%; whizzy paragraph -pdf xpdf -latex ./whizzypdfptex.sh
%; whizzy-paragraph "^\\\\begin{frame}"
% latex beamer presentation.
% platex, latex-beamer でコンパイルすることを想定。 

%     Tokyo Debian Meeting resources
%     Copyright (C) 2009 Junichi Uekawa
%     Copyright (C) 2009 Nobuhiro Iwamatsu

%     This program is free software; you can redistribute it and/or modify
%     it under the terms of the GNU General Public License as published by
%     the Free Software Foundation; either version 2 of the License, or
%     (at your option) any later version.

%     This program is distributed in the hope that it will be useful,
%     but WITHOUT ANY WARRANTY; without even the implied warreanty of
%     MERCHANTABILITY or FITNESS FOR A PARTICULAR PURPOSE.  See the
%     GNU General Public License for more details.

%     You should have received a copy of the GNU General Public License
%     along with this program; if not, write to the Free Software
%     Foundation, Inc., 51 Franklin St, Fifth Floor, Boston, MA  02110-1301 USA

\documentclass[cjk,dvipdfmx,12pt]{beamer}
\usetheme{Tokyo}
\usepackage{monthlypresentation}

%  preview (shell-command (concat "evince " (replace-regexp-in-string
%  "tex$" "pdf"(buffer-file-name)) "&")) 
%  presentation (shell-command (concat "xpdf -fullscreen " (replace-regexp-in-string "tex$" "pdf"(buffer-file-name)) "&"))
%  presentation (shell-command (concat "evince " (replace-regexp-in-string "tex$" "pdf"(buffer-file-name)) "&"))

%http://www.naney.org/diki/dk/hyperref.html
%日本語EUC系環境の時
\AtBeginDvi{\special{pdf:tounicode EUC-UCS2}}
%シフトJIS系環境の時
%\AtBeginDvi{\special{pdf:tounicode 90ms-RKSJ-UCS2}}

\title{Debian 9 Stretch のネットワークインターフェース名について}
\subtitle{Debian / Ubuntu ユーザーミートアップ in 札幌 2017.07}
\author{吉野 与志仁}
\date{2017年7月14日}
\logo{\includegraphics[width=8cm]{image200607/openlogo-light.eps}}

\begin{document}

\frame{\titlepage{}}

%\section{}

\begin{frame}{自己紹介}

\begin{itemize}
 \item 吉野 与志仁（よしの よしひと）
 \item 東京のほうから来ました
 \item @yy\_y\_ja\_jp
 \item Debian 公式開発者ではないです
 \item manpages-ja パッケージのメンテナ
 \item Debian JP Project メンバー
\end{itemize} 
\end{frame}

% \section{Agenda}
% \begin{frame}{Agenda}
%  Debian 9 Stretch のネットワークインターフェース名について
% \begin{enumerate}
%  \item 
% % XXX
% \end{enumerate}
% \end{frame}

\section{ネットワークインターフェース名}
\begin{frame}{ネットワークインターフェース名}
 デフォルトの名前が変わりました
 \pause
 
 \begin{itemize}[<+->]
  \item Debian 8 Jessieまで
	\begin{itemize}
	 \item 有線LANインターフェース: eth0, eth1, ...
	 \item 無線LANインターフェース: wlan0, wlan1, ...
	\end{itemize}
  \item Debian 9 Stretchから
	\begin{itemize}
	 \item 有線LANインターフェース: enp0s1 など、ハードウェア構成に
	       より異なる
	 \item 無線LANインターフェース: wlp1s0 など、ハードウェア構成に
	       より異なる
	\end{itemize}
 \end{itemize}
 
\end{frame}

%\section{アップグレード}
\begin{frame}{Jessie からのアップグレード}
 Jessie からアップグレードしたときは
 \pause

 \begin{itemize}[<+->]
  \item 今まで使っていたインターフェースの名前は変わりません
  \item 今後新たに使うインターフェースは新しい名前の形式になります -- 新
	しいUSB NICを差したときなど
 \end{itemize}
 \pause
 {\small\texttt{/usr/share/doc/udev/README.Debian.gz}}によると
 Debian 10 Busterでは昔の方式はサポートされないと言ってるので、新しい名前に移行したほうが
 よいかもしれません
\end{frame}

\begin{frame}{なぜ名前が変わったのか}
 今までの名前の変え方だとうまくいかないことがあったから

 \pause
 \vspace{2em}
 \texttt{eth0: renamed from eth1}
\end{frame}

\begin{frame}{なぜ名前を変えているのか}
 \begin{itemize}[<+->]
  \item Linux カーネルはネットワークインターフェースを認識するたびに名前を付
	けている eth0, eth1, ...
  \item 認識する順番は保証されない\\ \pause
	⇒ ネットワークインターフェースが複数あるときは、同じインター
	フェースが起動するたびに別の名前になることがある\\ \pause
	⇒ カーネルで認識されたらユーザーランド側で名前を修正することに
	した
 \end{itemize}
\end{frame}

\begin{frame}{名前の変え方 -- Jessieまで}
 udevパッケージの

 {\small{\texttt{/lib/udev/rules.d/75-persistent-net-generator.rules}}} と
 {\small{\texttt{/lib/udev/write\_net\_rules}}}\\
 \pause
 （Stretch にはもう存在しない！）

 \pause
 知らないインターフェースが現れたら
 \pause
 \begin{enumerate}[<+->]
  \item 空いている次の番号の名前を探す: eth1 など
  \item 「そのインターフェースが現れたらその名前に変える」というudevルールを\\
	{\small{\texttt{/etc/udev/rules.d/70-persistent-net.rules}}} に
	書き込む
  \item その名前に変える
 \end{enumerate}
 \pause
 Stretchでは無くなったため、
 \small{\texttt{/etc/udev/rules.d/70-persistent-net.rules}} が今後自動で書き換
 わることはありません
\end{frame}

\begin{frame}{名前の変え方 -- Stretchから}
 udevパッケージの
 {\footnotesize{\texttt{/lib/udev/rules.d/75-net-description.rules}}}、
 udevの{\footnotesize{\texttt{net\_id}}}ビルトイン、
 {\footnotesize{\texttt{/lib/udev/rules.d/80-net-setup-link.rules}}}、
 {\footnotesize{\texttt{net\_setup\_link}}}ビルトイン、
 {\footnotesize{\texttt{/lib/systemd/network/99-default.link}}} など

 \pause
 知らないインターフェースが現れたら
 \pause

 \begin{enumerate}[<+->]
  \item そのインターフェースを特定する情報（ハードウェア配置、MACアド
	レス、ユーザー設定ファイルなど）を取得する
  \item その情報に基づいて名前を決める: PCIバス0のスロット1に差さって
	いるとき enp0s1 など\\ \pause
	USB NICの場合は、USBの差す位置で名前が変わって
	ほしくないのでMACアドレスベースの名前にする設定になっている({\footnotesize{\texttt{/lib/udev/rules.d/73-usb-net-by-mac.rules}}}): enxaabbccxxyyzz など
 \end{enumerate}
\end{frame}

\begin{frame}{名前の変え方 -- （参考）一時期のUbuntu}
 Debianにはないが、Ubuntuはbiosdevnameパッケージで名前を決めている時期があった

 \begin{itemize}[<+->]
  \item インターフェース名: em0, p1p0 など
  \item 後発のudev net\_idはさらに別の名前形式にした
 \end{itemize}
\end{frame}

\begin{frame}[containsverbatim]{移行するには}

 \begin{enumerate}%[<+->]
  \item どんな名前になるか調べる eth0なら\\
	{\small{\texttt{udevadm test-builtin net\_id /sys/class/net/eth0}}}
       {\scriptsize
       \begin{verbatim}
# udevadm test-builtin net_id /sys/class/net/eth0
calling: test-builtin
（略）
Parsed configuration file /lib/systemd/network/99-default.link
Created link configuration context.
ID_NET_NAME_MAC=enx000d0bxxyyzz
ID_OUI_FROM_DATABASE=BUFFALO.INC
ID_NET_NAME_PATH=enp0s20u2
Unload module index
Unloaded link configuration context.
\end{verbatim}}
	{\footnotesize\texttt{ID\_NET\_NAME\_*}}のうち、順に
	{\footnotesize\texttt{ONBOARD}}、{\footnotesize\texttt{SLOT}}、
	{\footnotesize\texttt{PATH}}がもしあればそれが使われる({\footnotesize\texttt{/lib/systemd/network/99-default.link}})、ただしUSB NICは{\footnotesize\texttt{MAC}}が使われる
  \item 今までの名前を使っている設定ファイルを置き換える
  \item \small{\texttt{/etc/udev/rules.d/70-persistent-net.rules}} を消
	すかどこかに移動する
  \item 再起動
 \end{enumerate}
\end{frame}

\begin{frame}[containsverbatim]{移行せず、名前を変えさせないには}
 もともとeth0, wlan0しかなかったのに変えてほしくないときなど

 いくつかの方法
\begin{itemize}
 \item 空っぽの設定ファイルを\texttt{/etc}に置いて上書きする\\
       {\small{\verb|ln -s /dev/null /etc/systemd/network/99-default.link|}}
 \item カーネルコマンドライン引数に\texttt{net.ifnames=0}を追加\\
       /etc/default/grub に書いて update-grub を実行
\end{itemize}

 もちろんカーネルが決めた名前そのままになるので、インターフェースが複数
 あったらうまくいかないことが起きるでしょう
\end{frame}

\begin{frame}[containsverbatim]{名前を自分で設定するには}
\begin{itemize}
 \item \texttt{net\_setup\_link}のユーザー設定ファイル {\small{\texttt{/etc/systemd/network/*.link}}}
       を作る(systemd.link(5))\\
       例えばMACアドレス aa:bb:cc:xx:yy:zz のインターフェースを ethernet1 という名前にしたいなら
       {\footnotesize
       \begin{verbatim}
[Match]
MACAddress=aa:bb:cc:xx:yy:zz
[Link]
Name=ethernet1
\end{verbatim}} ただし、USB NICの名前を設定するには\\
       {\footnotesize{\verb|ln -s /dev/null /etc/udev/rules.d/73-usb-net-by-mac.rules|}}\\
       が必要
 \item udevルールファイル {\small{\texttt{/etc/udev/rules.d/*.rules}}} を作
       る(udev(7))\\
       Jessieまで自動で書き込まれてきた 
       {\small{\texttt{/etc/udev/rules.d/70-persistent-net.rules}}} に同
       じように追加するなど
\end{itemize}
\end{frame}

\begin{frame}[containsverbatim]{うまくいかないときは}
 udev をデバッグするしかない
\begin{itemize}
 \item 書いた\texttt{*.link}が期待通りに動くか見るには {\small{\texttt{net\_setup\_link}}}ビルトインを試す:\\
       {\footnotesize{\verb|udevadm test-builtin net_setup_link /sys/class/net/eth0|}}
 \item udevをdebugモードで起動する
       {\footnotesize
       \begin{verbatim}
# invoke-rc.d udev stop
# /lib/systemd/systemd-udevd --debug
\end{verbatim}}
 \item マシン起動時のudevをdebugモードにする\\
       /etc/udev/udev.conf を {\small{\verb|udev_log="debug"|}} にして
       {\small{\verb|dpkg-reconfigure linux-image-`uname -r`|}}
       して再起動
\end{itemize}

 名前変更に失敗してカーネルが決めた名前のままになっていることがある...
 見つけたらバグレポートしたほうがよいかもしれません
\end{frame}

\section{まとめ}
\begin{frame}[containsverbatim]{まとめ}
\begin{itemize}
 \item udevが変わり、Debian 9 Stretch からはネットワークインターフェースが新しい名前の
       形式になった
 \item 新しいインターフェースは新しい名前になる
 \item Jessie からアップグレードした場合は直ちに影響はないが、変えたほう
       がいいらしい
 \item 移行前に新しい名前を見るには\\
 {\footnotesize{\texttt{udevadm test-builtin net\_id /sys/class/net/<今の名前>}}}\\
 などが使える
 \item 名前付け替えを一切したくないときは \texttt{/etc} で
       \texttt{99-default.link}を上書き設定するとよい
 \item 自分で名前を付けたいときは
       {\small{\texttt{/etc/systemd/network/*.link}}} を作る、
       または昔ながらの udev ルール
       {\small{\texttt{/etc/udev/rules.d/*.rules}}} を作る
 \item うまく行かないときは udev をデバッグする。
       名前変更に失敗して元の名前のままになっていることがある
\end{itemize}
\end{frame}


\section{参考文献}
\begin{frame}[containsverbatim]{参考文献}
\begin{itemize}
 \item PredictableNetworkInterfaceNames
      {\scriptsize\verb|http://www.freedesktop.org/wiki/Software/systemd/PredictableNetworkInterfaceNames/|}
 \item systemd ソースパッケージ \verb|src/udev/udev-builtin-net_id.c|
 \item udev パッケージ \verb|/usr/share/doc/udev/README.Debian.gz|, man
       ページ \verb|systemd.link(5)|, \verb|udev(7)|
 \item biosdevname パッケージ(Ubuntu)
\end{itemize}
\end{frame}

%XXX
%

\end{document}

;;; Local Variables: ***
;;; outline-regexp: "\\([ 	]*\\\\\\(documentstyle\\|documentclass\\|emtext\\|section\\|begin{frame}\\)\\*?[ 	]*[[{]\\|[]+\\)" ***
;;; End: ***
